\begin{frame}{$\mu$ 항이 주인공, $\sigma$ 항은 보조 변수}
  \begin{itemize}
    \item 핵심 구조는 \kb{Policy Gradient Theorem의 one-vs-rest
      일반화}. 10.2절 softmax 그래디언트
      $[(1-p_a)s,\ -p_as]$ 의 ``고른 행동 확률 올리고, 나머지는
      내린다''가, 정규분포에서는 \kb{``고른 행동 쪽으로 평균을
      이동시킨다''}로 바뀜.
    \item \kb{$\mu$ 항}: 부호는 $a-\mu$ 의 부호로 정해짐.
      $A_t>0$ 이면 고른 행동 쪽으로 평균이 당겨지고, $A_t<0$
      이면 반대로 밀려남 — softmax의 ``고른 행동을 올린다''가
      정확히 이 $\mu$ 항.
    \item \kb{$\sigma$ 항}: 부호가 $(a-\mu)^2-\sigma^2$ 로 정해져
      까다로움.
      \begin{itemize}
        \item $|a-\mu|>\sigma$ (``의도하지 않은 극단적 행동'')
          $\to$ $\sigma$ \textbf{줄이는} 방향.
        \item $\sigma$ 이내(보통의 경우) $\to$ $\sigma$
          \textbf{늘리는} 방향.
      \end{itemize}
  \end{itemize}
  \begin{block}{기억하면 충분한 큰 그림}
    \kb{$\mu$ 가 주된 학습 대상이고, $\sigma$ 는 $A_t$ 의 부호와
    엔트로피 보너스 사이에 놓인 보조 변수}.
  \end{block}
\end{frame}
